
\documentclass[12pt]{article}
\usepackage[margin=1in]{geometry}
\usepackage{setspace}
\usepackage{parskip}

\title{Peer Review: Raja Chinnakotla --- Third Draft}
\author{Alejandro De La Torre}
\date{April 2026}

\begin{document}

\maketitle

\singlespacing

Before I get into the critique, I want to say: this is a genuinely ambitious project, and the fact that you pursued it at all is worth naming. Topics like this one get understudied (presumably) because the data is hostile — incomplete, politically contested, often buried in NGO archives rather than clean statistical series. A lot of people (me included) have walked away from questions like this because we knew the data wouldn't cooperate. You didn't, and that's real. Everything below is in the spirit of a paper worth taking seriously.

\noindent \textbf{Summary and Motivation.} \\
To my understanding, the paper asks how district-level demographic composition in Gujarat relates to the intensity of violence during the 2002 riots, using newly compiled data on the destruction of Muslim religious monuments alongside 2001 Census data. I found the question important and well-motivated, especially given how central these events are to the literature on inter-communal violence. I also appreciate the effort you put into assembling data from multiple sources—this is a real special sauce of the paper.

\vspace{0.5em}
\noindent \textbf{What I found most compelling.} \\
I think the strongest part of the paper is the attempt to bring together demographic composition (language and religion) with a concrete outcome that captures violence in a tangible way. The idea of focusing on monument destruction as a proxy for targeting is thoughtful, and the connection you draw to media exposure through language is an interesting angle that is not always the primary focus in this literature. I also found your motivation engaging, especially the effort to connect historical context to an empirical question.

\vspace{0.5em}
\noindent \textbf{Main suggestions for improvement.}

\textit{1. Clarifying the central contribution.} \\
I found myself wanting a clearer and more focused statement of what the paper is ultimately contributing. At the moment, the introduction and abstract point in a few different directions (language, religion, causal effects, and broader implications for the literature). I would suggest trying to distill the contribution into one concise sentence that clearly states what the paper shows and how it relates to existing work. Even reframing the paper more modestly as documenting district-level correlations, rather than making causal claims, could make the argument feel more precise and grounded.

\textit{2. Framing and interpretation of results.} \\
I think the discussion of the results could be tightened to better match what the empirical model can actually support. Since the estimates are not statistically significant and the sample size is small, I would be cautious about drawing broader conclusions about the literature. I found it more convincing to interpret the results as suggestive patterns rather than evidence for or against existing theories. A slightly more measured framing here would strengthen the credibility of the paper.

\textit{3. Data and outcome construction.} \\
I thought your choice of outcome variable was interesting, but I found myself wondering how it relates to underlying exposure. In particular, districts may differ in the baseline number of monuments, which could affect how we interpret destruction counts. I would suggest either discussing this more explicitly or considering a normalized measure (for example, destruction relative to baseline stock, if such data are available). Even a short discussion of this issue would help the reader better interpret the results.

\textit{4. Model and identification.} \\
I think it would help to clarify how you want the reader to interpret the regression model. I read the current specification as a cross-sectional association rather than a causal design, but the language in a few places suggests causal interpretation. I would suggest being explicit about this—either by reframing the results as correlations or by more clearly explaining what assumptions would be required for a causal interpretation. Making this distinction clear would go a long way in strengthening the paper.

\textit{5. Exposition and flow.} \\
I found the overall structure clear, but I think a few revisions could make the paper easier to follow, especially for readers without a background in South Asian studies. In particular, the research question could be stated more directly at the beginning of the introduction, and the transition from historical background to empirical analysis could be made more explicit. I also found that simplifying a few longer sentences would help the main ideas come through more clearly.

\vspace{0.5em}
\noindent \textbf{Overall.} \\
Overall, I think this is a thoughtful and promising paper with a strong topic and clear effort in data collection. My main suggestions focus on sharpening the contribution, aligning the claims more closely with the empirical design, and clarifying the presentation. With these adjustments, I think the paper will read as much more focused and persuasive while preserving its core idea.

\end{document}
